% Test of the minimum subset size k_min = 2 (table: reproducibility/figures_tables/kmin_tables.py;
% analyses: scripts/run_test_kmin.sh, analysis/compare_single_vs_pair.py).

\snote{kmin}

\emph{Question.} The method adds biological parameters until re-estimating them reproduces the full model and then removes those that are not needed, but the admissibility rule requires $|S|\ge2$, so the search cannot stop at one parameter and backward elimination cannot go below two. This is why 28 of the 31 admissible subsets have exactly two biological parameters: in those problems the size is set by the rule, not by the data, and in 15 of the 27 problems tested a single parameter already reaches the threshold (\sref{single}). This note asks whether requiring a second parameter has a measurable benefit or is an arbitrary choice.

\emph{Design.} For each of these 15 problems, the best single parameter $j^\ast$ (the admissible single parameter with the smallest median $e_{\mathrm{fit}}$ at $\pm$5\%, re-estimated together with the calibration parameters) was put through the same tests as the pair selected by our method: robustness to the size of the parameter changes, from $\pm$1\% to $\pm$50\% with 15 scenarios per size and with paired and with independent directions (\sref{robust}); admissibility at $\pm$20\% with the fixed scenario set; prediction of synthetic data with 20 seeds (\sref{synth}); and practical identifiability by the profile likelihood of the reduced problem (\sref{identif}). The rule $|S|\ge2$ is not applied to $\{j^\ast\}$. Pair and single parameter are compared problem by problem with the exact sign test (\texttt{scripts/run\_test\_kmin.sh}, \texttt{analysis/compare\_single\_vs\_pair.py}; 27~min of elapsed time on an 88-core server and 5.7~h of processor time, 4.7~h of it for Smith; \sref{time}).

\emph{Results} (Supplementary Table~\ref{tab:kmin}).
\begin{itemize}\itemsep0pt
\item \emph{The median criterion does not distinguish one parameter from two.} With one parameter, the median $e_{\mathrm{fit}}$ stays below $\tau_e$ at every perturbation size in 14 of the 15 problems (pattern A; with the pair, 15), and both are admissible at $\pm$20\% in all 15 problems. The exception is Borghans, whose single parameter fails at $\pm$1\% and, with independent directions, at several larger sizes.
\item \emph{Identifiability and prediction are the same.} Every selected parameter is identifiable in 12 of the 15 problems with the single parameter and in 10 with the pair ($P=0.5$): a single parameter cannot be compensated by another selected parameter. The synthetic success rate is 55\% (105 of 190 cases with signal) with the single parameter and 59\% (113 of 190) with the pair; the pair is better in 3 problems and the single parameter in 4 ($P=1$).
\item \emph{The pair fails in fewer individual scenarios and degrades less for large changes.} Over all perturbation sizes, 94\% of the scenarios reproduce the full model with the pair and 88\% with the single parameter with paired directions (the pair is better in 10 problems, the single parameter in 1; $P=0.012$), and 92\% against 87\% with independent directions (11 against 1; $P=0.006$). At $\pm$50\% with independent directions the median $e_{\mathrm{fit}}$ is 10.9\% with the pair and 18.7\% with the single parameter, larger with the single parameter in 14 of the 15 problems ($P<0.001$; with paired directions 10.5\% against 11.7\%, 13 of 15, $P=0.007$).
\end{itemize}

\emph{Interpretation.} A second parameter is not needed for the reduced model to be admissible in the median, identifiable or predictive: in about half of the problems a single biological parameter, re-estimated with the calibration parameters, already does all three. What the second parameter adds is a margin of safety: the reduction fails in fewer individual parameter changes and its error grows less when the parameters change by large amounts. The minimum size $k_{\min}=2$ is therefore a conservative choice, supported by this margin, rather than a requirement of the method; users who prefer the smallest possible model can set $k_{\min}=1$ (\texttt{config.py}), at the cost of a reduction that is less robust to individual large changes.

{\scriptsize\setlength{\tabcolsep}{2.5pt}
\begin{longtable}{@{}>{\raggedright\arraybackslash}p{1.7cm}>{\raggedright\arraybackslash}p{2.6cm}>{\raggedright\arraybackslash}p{3.0cm}ccccc@{}}
\caption{Best single parameter against the pair selected by our method, in the 15 problems in which one parameter already reaches the admissibility threshold at $\pm$5\%. Pass rate: percentage of the scenarios, over all perturbation sizes from $\pm$1\% to $\pm$50\%, with $e_{\mathrm{fit}}\le\tau_e$, with paired (P) and independent (I) directions; $e_{\mathrm{fit}}$ at $\pm$50\%: median over the 15 scenarios with independent directions (\%); $\pm$20\%: admissible at $\pm$20\%; Synthetic: cases with signal in which $e_{\mathrm{pred}}\le\tau_e$ over the 20 seeds (identifiability is reported in the text). Each cell gives pair / single.}\label{tab:kmin}\\
\toprule
Problem & Single parameter & Pair & Pass rate P (\%) & Pass rate I (\%) & $e_{\mathrm{fit}}$ $\pm$50\% & $\pm$20\% & Synthetic \\
\midrule
\endfirsthead
\toprule
Problem & Single parameter & Pair & Pass rate P (\%) & Pass rate I (\%) & $e_{\mathrm{fit}}$ $\pm$50\% & $\pm$20\% & Synthetic \\
\midrule
\endhead
Alkan & \texttt{kt} & \texttt{kt}, \texttt{k\_\allowbreak{}rep\_\allowbreak{}nhej} & 100 / 81 & 96 / 83 & 17 / 26 & yes / yes & 3/3 / 3/3 \\
Borghans & \texttt{v0} & \texttt{Kz}, \texttt{K\_\allowbreak{}par} & 78 / 68 & 66 / 51 & 43 / 47 & yes / yes & 2/11 / 1/11 \\
Brannmark & \texttt{km2} & \texttt{k1f}, \texttt{km2} & 93 / 91 & 99 / 99 & 13 / 19 & yes / yes & 1/18 / 1/18 \\
Fiedler & \texttt{k10} & \texttt{tau2}, \texttt{k6} & 100 / 100 & 99 / 99 & 12 / 13 & yes / yes & 1/6 / 1/6 \\
Fujita & \texttt{reaction\_\allowbreak{}2\_\allowbreak{}k1} & \texttt{init\_\allowbreak{}AKT}, \texttt{reaction\_\allowbreak{}6\_\allowbreak{}k1} & 72 / 93 & 70 / 86 & 27 / 15 & yes / yes & 14/18 / 16/18 \\
Giordano & \texttt{gamma\_\allowbreak{}4} & \texttt{alpha\_\allowbreak{}4}, \texttt{alpha\_\allowbreak{}28} & 92 / 80 & 97 / 86 & 6 / 33 & yes / yes & -- / -- \\
Laske & \texttt{k\_\allowbreak{}imp} & \texttt{k\_\allowbreak{}imp}, \texttt{k\_\allowbreak{}syn\_\allowbreak{}R\_\allowbreak{}M} & 100 / 80 & 98 / 80 & 11 / 27 & yes / yes & -- / -- \\
Okuonghae & \texttt{gamma\_\allowbreak{}0} & \texttt{transmission\_\allowbreak{}rate\_\allowbreak{}effective}, \texttt{gamma\_\allowbreak{}i} & 99 / 99 & 96 / 94 & 6 / 7 & yes / yes & 17/20 / 18/20 \\
Oliveira & \texttt{delta\_\allowbreak{}} & \texttt{beta\_\allowbreak{}0}, \texttt{h\_\allowbreak{}hosp\_\allowbreak{}rate} & 94 / 91 & 90 / 89 & 5 / 17 & yes / yes & 12/20 / 12/20 \\
Rahman & \texttt{infected\_\allowbreak{}moderate\_\allowbreak{}transmission\_\allowbreak{}rate} & \texttt{infected\_\allowbreak{}moderate\_\allowbreak{}transmission\_\allowbreak{}rate}, \texttt{infected\_\allowbreak{}normal\_\allowbreak{}worsen\_\allowbreak{}rate} & 95 / 84 & 99 / 83 & 3 / 22 & yes / yes & 15/19 / 17/19 \\
Raia & \texttt{init\_\allowbreak{}Rec\_\allowbreak{}i} & \texttt{SOCS3mRNA\_\allowbreak{}production}, \texttt{init\_\allowbreak{}Rec\_\allowbreak{}i} & 100 / 97 & 98 / 94 & 3 / 5 & yes / yes & 7/8 / 7/8 \\
Salazar-Cavazos & \texttt{ratio\_\allowbreak{}kpkd\_\allowbreak{}Y1068\_\allowbreak{}\_\allowbreak{}FREE} & \texttt{ratio\_\allowbreak{}kpkd\_\allowbreak{}Y1068\_\allowbreak{}\_\allowbreak{}FREE}, \texttt{ratio\_\allowbreak{}kpkd\_\allowbreak{}YN\_\allowbreak{}\_\allowbreak{}FREE} & 87 / 67 & 78 / 62 & 15 / 28 & yes / yes & 12/18 / 3/18 \\
Schwen & \texttt{ini\_\allowbreak{}R2fold} & \texttt{ini\_\allowbreak{}R1}, \texttt{ini\_\allowbreak{}R2fold} & 93 / 89 & 97 / 94 & 15 / 19 & yes / yes & 0/9 / 0/9 \\
Smith & \texttt{k30r} & \texttt{k8}, \texttt{kminus7b} & 100 / 100 & 100 / 99 & 8 / 10 & yes / yes & 18/20 / 19/20 \\
Weber & \texttt{a21} & \texttt{s12}, \texttt{a21} & 100 / 100 & 100 / 100 & 6 / 6 & yes / yes & 11/20 / 7/20 \\
\bottomrule
\end{longtable}

}
